Fix an integer $r\ge 2$.
For a graph $G$ and integer $m\ge 1$, consider the property:

(P$_m$): Every (not necessarily induced) subgraph $H\subseteq G$ on $m$ vertices can be written as a union $H=H_0\cup H_1$ where $\chi(H_0)\le r$ and $|E(H_1)|\le t$.

Define $f_r(m)$ to be the largest integer $t$ (as a function of $m$) such that the following implication holds for all graphs $G$:

If $G$ satisfies (P$_m$) with this $t=f_r(m)$ for every $m$ (i.e. for each $m$-vertex subgraph $H$ we can decompose with at most $f_r(m)$ edges in $H_1$), then $\chi(G)\le r+1$.

The question asks whether $f_2(m)\gg m$ (linear growth), and more generally whether $f_r(m)\gg_r m$.
